\section*{Bab \ref{ch:g7:areas} --- Luas dan Volume}

\begin{solution}{exo:g7:areas:1}
$8 \times 5 = 40$ cm$^2$; \qquad $6.5 \times 4 = 26$ m$^2$.
\end{solution}

\begin{solution}{exo:g7:areas:2}
$A = 10 \times 4 = 40$ cm$^2$. Jawabannya bukan
$10 \times 6 = 60$ karena sisi $6$ cm itu miring: rumusnya
memakai \emph{tinggi} (jarak \omterm{def:g6:lines:perp}{tegak lurus}), yang $4$ cm.
\end{solution}

\begin{solution}{exo:g7:areas:3}
$\frac{12 \times 7}{2} = 42$ cm$^2$; \quad
$\frac{9 \times 6}{2} = 27$; \quad
$\frac{5.5 \times 2}{2} = 5.5$ m$^2$.
\end{solution}

\begin{solution}{exo:g7:areas:4}
Cakram: $\pi \times 5^2 = 25\pi \approx 79$ cm$^2$. Setengah cakram
berjari-jari $10$: $\frac{\pi \times 10^2}{2} = 50\pi \approx 157$ cm$^2$.
\end{solution}

\begin{solution}{exo:g7:areas:5}
\omterm{def:g6:measure:perimeter}{Kelilingnya}: $2\pi \times 3 = 6\pi \approx 18.8$ cm (sebuah panjang,
dalam cm). Luasnya: $\pi \times 3^2 = 9\pi \approx 28.3$ cm$^2$
(sebuah permukaan, dalam cm$^2$).
\end{solution}

\begin{solution}{exo:g7:areas:6}
\omterm{def:g7:areas:prism}{Prisma}: $24 \times 10 = 240$ cm$^3$. \omterm{def:g7:areas:prism}{Tabung}:
$\pi \times 3^2 \times 8 = 72\pi \approx 226$ cm$^3$.
\end{solution}

\begin{solution}{exo:g7:areas:7}
$\frac{8 \times h}{2} = 28$, jadi $4h = 28$ dan $h = 7$ cm.
\end{solution}

\begin{solution}{exo:g7:areas:8}
Ketiga \omterm{def:g2:mult:def}{hasil kalinya} sesuai karena masing-masing menghitung besaran
yang \emph{sama} --- yaitu luas segitiganya. Ini memberi pemeriksaan
yang praktis untuk pelukisan, dan cara menghitung sebuah tinggi dari
pasangan alas--tinggi yang lain (dipakai pada \cref{exo:g7:areas:12}).
\end{solution}

\begin{solution}{exo:g7:areas:9}
Kedua segitiganya punya luas $\frac{30 \times 20}{2} = 300$ m$^2$
dan $\frac{18 \times 20}{2} = 180$ m$^2$: seluruhnya $480$ m$^2$.
Rumus umum yang disarankan: untuk sisi \omterm{def:g6:lines:perp}{sejajar} $a$ dan $b$ pada jarak $h$,
\[
A = \frac{(a + b) \times h}{2}
\]
--- yaitu rumus trapesium (di sini $\frac{(30+18)\times 20}{2} = 480$).
\end{solution}

\begin{solution}{exo:g7:areas:10}
\omterm{def:g6:measure:volume}{Volume} airnya: $\pi \times 6^2 \times 15 = 540\pi \approx 1\,696$
cm$^3$. Alas \omterm{def:g5:solids:def}{baloknya}: $18 \times 12 = 216$ cm$^2$. Tinggi yang
tercapai: $\frac{540\pi}{216} = 2.5\pi \approx 7.9$ cm.
\end{solution}

\begin{solution}{exo:g7:areas:11}
\omterm{def:g3:shapes:shapes}{Jari-jarinya} $15$ dan $20$ cm. Luasnya: $225\pi \approx 707$ cm$^2$
dan $400\pi \approx 1257$ cm$^2$. Harga per $100$ cm$^2$:
$\frac{9}{7.07} \approx 1.27$ euro dan
$\frac{14}{12.57} \approx 1.11$ euro. Piza yang besar lebih
menguntungkan --- luas tumbuh dengan \emph{kuadrat} \omterm{def:g4:geometry:circle}{diameternya}.
\end{solution}

\begin{solution}{exo:g7:areas:12}
Dengan sisi siku-sikunya: $A = \frac{6 \times 8}{2} = 24$ cm$^2$.
Dengan sisi miringnya sebagai alas: $24 = \frac{10 \times h}{2} = 5h$,
jadi $h = 4.8$ cm.
\end{solution}

\begin{solution}{pb:g7:areas:1}
\textbf{1.} Persegi: $8 \times 8 = 64$. Persegi panjang:
$5 \times 13 = 65$. Keempat kepingan yang sama tampak menutupi $64$
persegi satuan pada bangun yang satu dan $65$ pada yang lain: satu
satuan luas tercipta dari ketiadaan --- katanya begitu.

\textbf{2.} Tiap segitiganya: $\frac{8 \times 3}{2} = 12$. Tiap
trapesiumnya: $\frac{(5 + 3) \times 5}{2} = 20$.

\textbf{3.} $12 + 12 + 20 + 20 = 64$: kepingannya berjumlah $64$,
jadi kepingannya dapat mengisi perseginya dengan tepat --- dan
\emph{tidak} dapat mengisi persegi panjang yang $65$ satuan. Gambar
persegi panjangnya pasti memuat sebuah lubang.

\textbf{4.} Kecuraman yang sama berarti naik $3$ sepanjang $8$
\omterm{def:g6:propdata:def}{sebanding} dengan naik $2$ sepanjang $5$; \omterm{def:g2:mult:def}{hasil kali} silangnya:
$3 \times 5 = 15$ dan $2 \times 8 = 16$. Tidak sama: sisi
miringnya dan tepi miringnya \emph{tidak} segaris.
``Diagonal'' persegi panjangnya adalah dusta --- dua kemiringan yang
sedikit berbeda yang bertemu pada sudut yang tumpul.

\textbf{5.} Di sepanjang diagonal palsunya keempat kepingan itu
meninggalkan celah yang panjang dan luar biasa tipis --- sebilah
tipis berbentuk \omterm{def:g7:central:parallelogram}{jajargenjang} yang sangat datar --- yang luasnya
persis satu persegi satuan yang hilang, $65 - 64 = 1$. Kalau
dibentangkan sepanjang diagonal yang kira-kira $13$ satuan
panjangnya, lebar rata-ratanya kira-kira
$1 \div 13 \approx 0.08$ satuan: pada gambar sungguhan, lebih tipis
daripada garis pensil yang menyembunyikannya.

\textbf{6.} Tiap segitiganya punya alas $6$ cm dan tinggi $4$ cm ---
yaitu jarak antara kedua \omterm{def:g6:lines:perp}{garis sejajarnya}, di mana pun puncaknya
duduk --- jadi tiap luasnya $\frac{6 \times 4}{2} = 12$ cm$^2$
(\cref{thm:g7:areas:triangle}). Menggeser puncaknya di sepanjang
\omterm{def:g6:lines:perp}{garis sejajarnya} mengubah bentuknya, tidak pernah alas atau
tingginya: luasnya sama selamanya.

\textbf{7.} Diagonalnya memotong trapesiumnya menjadi segitiga
beralas $B$ dan segitiga beralas $b$, keduanya setinggi $h$ (yaitu
jarak antara sisi \omterm{def:g6:lines:perp}{sejajarnya}):
\[
A = \frac{B \times h}{2} + \frac{b \times h}{2}
= \frac{(B + b) \times h}{2} .
\]

\textbf{8.} $A = \frac{(9 + 5) \times 4}{2} = \frac{56}{2} =
28$ cm$^2$.

\textbf{9.} Rata-rata sisi \omterm{def:g6:lines:perp}{sejajarnya} adalah
$\frac{B + b}{2}$, jadi (rata-rata) $\times h =
\frac{(B+b) \times h}{2}$: rumus yang sama, yang difaktorkan
secara berbeda. Periksa: rata-ratanya $= \frac{9+5}{2} = 7$, dan
$7 \times 4 = 28$ cm$^2$.

\textbf{10.} Tumpukan yang condong memuat kartu yang persis sama
--- lapisan yang sama, masing-masing dengan luas yang sama --- jadi
\omterm{def:g6:measure:volume}{volumenya} sama. Pada $V = B \times h$, tingginya harus diukur
\emph{\omterm{def:g6:lines:perp}{tegak lurus}} terhadap alasnya (yaitu tinggi tegak tumpukannya),
bukan sepanjang condongnya: persis seperti luas \omterm{def:g7:central:parallelogram}{jajargenjang} yang
menghendaki tinggi tegak lurusnya, bukan sisi miringnya
(\cref{exo:g7:areas:2}), satu dimensi lebih tinggi.

\textbf{11.} Cincinnya $= \pi \times 5^2 - \pi \times 3^2 = 25\pi -
9\pi = 16\pi \approx 50$ m$^2$ (\cref{thm:g7:areas:disk}).

\textbf{12.} Cakram dengan luas $16\pi$ berjari-jari $4$, karena
$\pi \times 4^2 = 16\pi$. \omterm{def:g3:shapes:shapes}{Jari-jarinya} $3$, $4$, $5$: tripel yang
paling terkenal dalam matematika, yang berperan utama pada
\cref{ch:g8:pythagoras}.

\textbf{13.} Satu piza $40$ cm: \omterm{def:g3:shapes:shapes}{jari-jarinya} $20$,
luasnya $\pi \times 400 \approx 1\,257$ cm$^2$. Dua piza $28$ cm:
\omterm{def:g3:shapes:shapes}{jari-jarinya} $14$, bersama-sama
$2 \times \pi \times 196 = 392\pi \approx 1\,232$ cm$^2$. Satu piza
yang besar itu \emph{lebih} banyak daripada dua piza sedang.

\textbf{14.} Menskalakan \omterm{def:g3:shapes:shapes}{jari-jarinya} dengan $1.4$ menskalakan
luasnya dengan $1.4^2 = 1.96$ --- hampir dua kali lipat. Faktor yang
persis adalah bilangan yang kuadratnya $2$, kira-kira $1.414$: yaitu
bilangan diagonal pada \cref{ex:g8:pythagoras:diagonal}. (Patokan
kasar: setengah kali lagi lebarnya berarti dua kali pizanya, hampir.)

\textbf{15.} $13 \times 13 = 169$, tetapi $8 \times 21 = 168$:
kali ini satu persegi satuan \emph{menghilang} --- kepingannya
bertindih pada sebilah tipis, bukannya meninggalkan celah. Dengan
suku berurutan $3, 5, 8, 13, 21$ pada barisan
Hemachandra--Fibonacci, \omterm{def:g2:mult:def}{hasil kalinya} berganti-ganti antara satu
lebih banyak dan satu lebih sedikit daripada kuadratnya
($5 \times 13 = 65 = 8^2 + 1$, lalu $8 \times 21 = 168 = 13^2 - 1$):
pesulapnya bergantian ``untung'' dan ``rugi'' satu satuan. Luas tidak
pernah berbohong: kepingan yang berjumlah $64$ menutupi $64$, ke mana
pun kepingannya digeser --- setiap satuan yang hilang atau berlebih
sedang bersembunyi pada sebilah tipis, menunggu pemeriksaan
kemiringan untuk menyingkapkannya.

\end{solution}
